\section{Formal Proofs}
\label{sec:proofs}

The results of \cref{sec:conservation} are not informal. Each is a
theorem in \texttt{Ratio.Core}, a dependency-free Lean~4~\cite{demoura2021lean}
development that
uses only the core library and the \texttt{omega} decision procedure for
linear integer arithmetic --- no Mathlib, no axioms beyond Lean's
kernel. This section gives the proofs at the level of detail a reader
needs to check them by hand; the Lean source is the machine-checked
source of record, and each statement is tagged with the name that
discharges it.

\subsection{The algebra is a commutative group}

The pointwise operations of \cref{def:vec} inherit their laws from
$\mathbb{Z}$ dimension by dimension.

\begin{lemma}[Group laws, pointwise]
\label{lem:group}
For all $\mathbf{a},\mathbf{b},\mathbf{c} \in \mathbb{Z}^{\rdim}$:
\[
\begin{aligned}
  \mathbf{a}+\mathbf{b} &= \mathbf{b}+\mathbf{a}, &
  (\mathbf{a}+\mathbf{b})+\mathbf{c} &= \mathbf{a}+(\mathbf{b}+\mathbf{c}), \\
  \mathbf{a}+\rzero &= \mathbf{a}, &
  \mathbf{a}+(-\mathbf{a}) &= \rzero.
\end{aligned}
\]
\leanref{Ratio.Vec.add\_comm/assoc/zero/neg}
\end{lemma}

\begin{proof}
Each equation is between functions $\rdim \to \mathbb{Z}$, so by
function extensionality it suffices to fix an arbitrary $d$ and prove the
scalar equation. After unfolding the pointwise definitions, every case
reduces to an identity of linear integer arithmetic --- e.g.\
$\mathbf{a}(d) + \mathbf{b}(d) = \mathbf{b}(d) + \mathbf{a}(d)$ for
commutativity, $\mathbf{a}(d) + (-\mathbf{a}(d)) = 0$ for the inverse ---
each closed by \texttt{omega}. In Lean the four proofs are one line apiece
(\texttt{funext d; show \ldots; omega}).
\end{proof}

\subsection{Conservation is a subgroup}

\begin{lemma}[Closure of conservation]
\label{lem:closure}
$\rconserves(\rzero)$; and for all $\mathbf{a},\mathbf{b}$,
\[
  \rconserves(\mathbf{a}) \wedge \rconserves(\mathbf{b})
  \Rightarrow \rconserves(\mathbf{a}+\mathbf{b}),
  \qquad
  \rconserves(\mathbf{a}) \Rightarrow \rconserves(-\mathbf{a}).
\]
\leanref{Ratio.conserves\_zero/add/neg}
\end{lemma}

\begin{proof}
$\rconserves(\rvec)$ unfolds to $\rvec = \rzero$, so $\rconserves(\rzero)$ holds
by reflexivity. For closure under $+$: from $\mathbf{a} = \rzero$ and
$\mathbf{b} = \rzero$ rewrite to get $\mathbf{a}+\mathbf{b} = \rzero + \rzero$,
which is $\rzero$ by \cref{lem:group}'s identity law. For closure under
negation: from $\mathbf{a} = \rzero$, $-\mathbf{a} = -\rzero = \rzero$
pointwise (\texttt{omega} on $-(0)=0$).
\end{proof}

\Cref{lem:closure} is exactly the statement that the balanced
transactions form a subgroup of $\mathbb{Z}^{\rdim}$. The conservation law
is then an induction that this subgroup absorbs the fold.

\subsection{The conservation law}

\begin{lemma}[Fold stays conserving]
\label{lem:fold}
For every log and every accumulator $\textit{acc}$, if
$\rconserves(\textit{acc})$ and every transaction in the log conserves,
then $\rconserves(\mathrm{foldl}\,(+)\,\textit{acc}\,\textit{log})$.
\leanref{Ratio.foldl\_conserves}
\end{lemma}

\begin{proof}
By induction on the log.
\emph{Base.} For the empty log the fold returns $\textit{acc}$, which
conserves by hypothesis.
\emph{Step.} For $\textit{log} = t :: ts$, the fold takes one step to
$\mathrm{foldl}\,(+)\,(\textit{acc}+t)\,ts$. The new accumulator
$\textit{acc}+t$ conserves by closure under $+$ (\cref{lem:closure}),
since $\textit{acc}$ conserves and $t \in (t::ts)$ conserves by
hypothesis. The remaining transactions $ts$ all conserve (they are a
sub-list), so the induction hypothesis applies and yields the result.
\end{proof}

\begin{theorem}[Conservation law]
\label{thm:conservation}
If every transaction in an append-only log conserves value, then the
folded ledger conserves value:
\[
  \big(\forall t \in \textit{log},\; \rconserves(t)\big)
  \;\Longrightarrow\;
  \rconserves\!\big(\rledger(\textit{log})\big).
\]
\leanref{Ratio.ledger\_conserves}
\end{theorem}

\begin{proof}
$\rledger(\textit{log}) = \mathrm{foldl}\,(+)\,\rzero\,\textit{log}$ by
\cref{def:ledger}. Apply \cref{lem:fold} with $\textit{acc} = \rzero$,
whose conservation is $\rconserves(\rzero)$ from \cref{lem:closure}.
\end{proof}

\Cref{thm:conservation} is the whole game. It is a closed-form guarantee,
quantified over \emph{all} logs of \emph{all} lengths, that a ledger
assembled from balanced transactions is itself balanced --- proved
structurally, checked by machine. No production accounting system we know
of can state its central invariant this way, let alone prove it.

\subsection{Double-entry, recovered}

To show the generalization loses nothing, we recover classical
double-entry as the one-dimensional instance. Let
$\mathrm{total}[a_1,\dots,a_n] = \sum_i a_i$ and call a split list
\emph{balanced} when its total is $0$.

\begin{theorem}[Two-leg postings balance]
\label{thm:double-entry}
For every amount $a$, the posting $[a, -a]$ is balanced:
$\mathrm{total}[a,-a] = 0$.
\leanref{Ratio.balanced\_pair}
\end{theorem}

\begin{proof}
$\mathrm{total}[a,-a] = a + (-a + 0) = 0$ by \texttt{omega}. The empty
posting is balanced by reflexivity (\texttt{balanced\_nil}). These
are the $|\rdim|=1$ shadows of \cref{lem:group}'s inverse and identity
laws.
\end{proof}

\subsection{What is and is not proved}

The kernel's \emph{arithmetic} is proved: the group structure, the
subgroup closures, and the conservation law hold for all inputs. What is
deliberately \emph{not} a theorem is any business claim --- that a
particular posting rule maps a particular event correctly, or that an
account's normal balance is debit. Those are policy, and policy lives in
the control plane (\cref{sec:control}). The normal-balance table
(\texttt{Ratio.normalSide}) is carried in Lean as a \emph{checked example}
of such policy --- assets and expenses are debit-normal, liabilities,
equity, and income credit-normal (\texttt{asset\_debit\_normal}, etc.) ---
precisely to demonstrate the boundary: it is decidable and checked, but it
is not a kernel invariant and the kernel does not depend on it.
\Cref{tab:trust} summarizes the trust boundary.

% The trust boundary as a table: what is proved vs. what is policy.
\begin{table}[t]
\centering
\renewcommand{\arraystretch}{1.25}
\begin{tabular}{@{}p{0.40\linewidth} l p{0.34\linewidth}@{}}
\toprule
\textbf{Concern} & \textbf{Status} & \textbf{Source of record} \\
\midrule
Group laws of $\mathbb{Z}^{\rdim}$            & proved   & \texttt{Vec.add\_comm/assoc/zero/neg} \\
Conservation is a subgroup                   & proved   & \texttt{conserves\_zero/add/neg} \\
Ledger fold preserves conservation           & proved   & \texttt{ledger\_conserves} \\
Double-entry (two-leg) balances              & proved   & \texttt{balanced\_pair} \\
Running \texttt{Money} arithmetic            & emitted  & \texttt{Ratio.Common.Emit} (from proofs) \\
\midrule
Normal-balance classification                & checked policy & \texttt{normalSide} (example, not invariant) \\
Posting rules / account mappings             & policy   & content-addressed control plane \\
Prices / FX / reference data                 & policy   & typed facts with provenance \\
NL intent $\to$ specification                 & untrusted & LLM proposal, compiled + checked \\
\bottomrule
\end{tabular}
\caption{The trust boundary. The top block is machine-checked (or emitted
from machine-checked sources) and forms the trusted computing base; the
bottom block is expressive policy that lives outside the kernel and cannot
violate conservation, because admission runs through
\cref{thm:conservation}.}
\label{tab:trust}
\end{table}

